\section{Per-constructor equations versus guarded disjunctive case axioms}
\label{sec:casesplit}

The translation of \cref{sec:translation} renders a pattern-matching
definition as per-constructor equations.  The translation of
\cite[\S 5.2]{CzajkaKaliszyk2018} renders the same definition as a single
\emph{guarded disjunctive axiom}.  This section compares the two forms
inside the framework just built: the model $\ML$ arbitrates which guards
are semantically necessary, and first-order derivations settle the
logical relationship.  Fix a definition by a case split on a variable,
\[
  c \defeq \lambda x{:}I.\ \tcase{x}{z.\tau}{\{\wcns{c}_i(\tup{y}_i)
  \Rightarrow b_i\}_i} \;:\; \Pi x{:}I.\,\tau ,
\]
and abbreviate $R_i \defeq \compf(\erase{b_i})$ (over the variables
$\tup{y}_i$).  Our translation emits the \emph{split form}
\begin{equation}
  \tag{$\mathrm{Eq}_c$}
  \forall \tup{y}_i.\ \Ihat{c}(\Ihat{\wcns{c}}_i(\tup{y}_i)) \feq R_i
  \qquad \text{(one axiom per constructor $i$),}
\end{equation}
whereas the disjunctive rendering is
\begin{equation}
  \tag{$\mathrm{Disj}_c$}
  \forall x.\ \Ihat{I}(x) \to \bigvee_i \exists \tup{y}_i.\
    {\textstyle\bigwedge_j} \guard{\sigma_{ij}}{y_{ij}} \wedge
    x \feq \Ihat{\wcns{c}}_i(\tup{y}_i) \wedge \Ihat{c}(x) \feq R_i ,
\end{equation}
and we write $\mathrm{Disj}^-_c$ for the variant with the guard
$\Ihat{I}(x)$ omitted.

\subsection{Semantic status}
\label{sec:casesplit-semantic}

\begin{proposition}[The split form needs no guards]
\label{prop:split-valid}
$\ML \vDash \mathrm{Eq}_c$, with the universally quantified variables
$\tup{y}_i$ ranging over the \emph{entire} domain $\Dom$ --- junk
included.  Moreover $\ML \vDash \mathrm{Disj}_c$.
\end{proposition}

\begin{proof}
The first claim is an instance of \cref{lem:compile-sound}: the emitted
equations are valid in $\ML$ under \emph{every} assignment; concretely,
for arbitrary $\tup{d} \in \Dom$,
$\Ihat{c}^{\ML}(\Ihat{\wcns{c}}_i^{\ML}(\tup{d})) =
\cls{c\ \wcns{c}_i(\tup{e})} \lconv \cls{\erase{b_i}^{[\tup{y}_i \mapsto
\tup{d}]}}$ by a $\delta$-, $\beta$- and $\iota$-step --- the reduction
fires syntactically, with no typing side condition.  For the second
claim: given $d \in \Ihat{I}^{\ML}$, inversion validity
(\cref{thm:env-validity}) yields a disjunct $i$ and witnesses with $d =
\Ihat{\wcns{c}}_i^{\ML}(\tup{e})$, and $\mathrm{Eq}_c$ supplies
$\Ihat{c}^{\ML}(d) = $ the corresponding $R_i$-value.
\end{proof}

\begin{proposition}[The disjunctive form's guard is mandatory]
\label{prop:disj-guard}
Let $\Env$ contain $\nat$ and the definition $\mathsf{mypred}$ of
\cref{ex:compile}.  Then $\ML \nvDash \mathrm{Disj}^-_{\mathsf{mypred}}$.
Moreover, guard omission in the disjunctive form is not merely false in
$\ML$ but \emph{absolutely} unsound: for a single-constructor datatype
$U$ with $\wcns{u} : ()$ and any definition $c_0$ by a (trivial) case on
$U$, the theory $\{\mathrm{Disj}^-_{c_0}\} \cup \{\text{discrimination
for } \bool\}$ is inconsistent.
\end{proposition}

\begin{proof}
For the first claim, take the junk element $d = \cls{\csO\,\csO} \in
\Dom$ (the constructor constant $\csO$ applied to itself).  By
\cref{cor:discrimination}(2), $d$ is not convertible to any constructor
form, so both disjuncts of $\mathrm{Disj}^-_{\mathsf{mypred}}$ fail at $x
\defeq d$: $d \neq \cls{\csO}$ and $d \neq \cls{\csS(e)}$ for every $e$.
(The guarded form is unaffected: $d \notin \semr{\nat}{} =
\Ihat{\nat}^{\ML}$, again by \cref{cor:discrimination}(2) and stage
induction.)

For the second claim, $\mathrm{Disj}^-_{c_0}$ has the single disjunct
$\forall x.\ x \feq \Ihat{\wcns{u}} \wedge \cdots$, which entails
$\forall x\,y.\ x \feq y$; instantiating with
$\Ihat{\cstrue}$ and $\Ihat{\csfalse}$ contradicts discrimination.
This reproduces, inside our framework, the counterexample that makes the
guard on the matched term's variables mandatory in
\cite[\S 5.6]{CzajkaKaliszyk2018}.
\end{proof}

The two propositions isolate the semantic content of the design change.
The disjunctive axiom asserts \emph{exhaustiveness}: every domain element
(in the guarded version, every element satisfying the guard) is of
constructor form.  Exhaustiveness is a property of the \emph{type}, is
false at junk, and therefore requires the guard.  The equations assert
only \emph{computation}: they are silent off their patterns, hold at
every point of the model, and need no guard --- neither on the matched
term nor on the constructor-argument variables $\tup{y}_i$.  In the
factored architecture this is no accident: exhaustiveness belongs to
specification extraction and is emitted once per datatype (the inversion
axiom, guarded), while the program's equations carry no typing content at
all.  The guard class that \cite[\S 5.6]{CzajkaKaliszyk2018} identifies
as impossible to omit is thereby not omitted but \emph{dissolved}.

\begin{remark}[Guard policy for constructor arguments]
\label{rem:guard-policy}
In the split form, guards on the fresh constructor-argument variables
$\tup{y}_i$ would also be sound (weakening a valid formula), and
correspond to the optional closure-variable guards of lambda-lifting
(\texttt{opt\_closure\_guards} in the implementation); by
\cref{prop:split-valid} they are unnecessary, uniformly with the
treatment of lambda-lifted closure variables --- but here this is a
theorem about $\ML$ rather than a conjecture.
\end{remark}

\subsection{Logical relationship}
\label{sec:casesplit-derive}

\begin{proposition}[Interderivability relative to structural axioms]
\label{prop:interderivable}
Let $\mathrm{Str}_I$ denote the structural axioms of $I$ emitted by
\cref{def:theory} (constructor typing, inversion, injectivity,
discrimination).  Then, in classical first-order logic with equality:
\begin{enumerate}[leftmargin=2em]
\item $\mathrm{Eq}_c \cup \{\text{inversion}\} \provesfol
  \mathrm{Disj}_c$;
\item $\{\mathrm{Disj}_c\} \cup \mathrm{Str}_I \provesfol$ the
  \emph{guarded} equations
  $\forall \tup{y}_i.\ \bigwedge_j \guard{\sigma_{ij}}{y_{ij}} \to
  \Ihat{c}(\Ihat{\wcns{c}}_i(\tup{y}_i)) \feq R_i$ \ for each $i$.
\end{enumerate}
The unguarded equations $\mathrm{Eq}_c$ are strictly stronger than
$\{\mathrm{Disj}_c\} \cup \mathrm{Str}_I$: they are not derivable from
it.
\end{proposition}

\begin{proof}
(1) Assume $\Ihat{I}(x)$.  The inversion axiom yields a disjunct $i$ with
witnesses $\tup{y}_i$, their guards, and $x \feq
\Ihat{\wcns{c}}_i(\tup{y}_i)$; rewriting with this equation in
$\Ihat{c}(x)$ and applying the $i$-th equation of $\mathrm{Eq}_c$
produces $\Ihat{c}(x) \feq R_i$, completing the $i$-th disjunct of
$\mathrm{Disj}_c$.

(2) Fix $i$ and assume the guards of $\tup{y}_i$.  By constructor
typing, $\Ihat{I}(\Ihat{\wcns{c}}_i(\tup{y}_i))$; instantiate
$\mathrm{Disj}_c$ at $x \defeq \Ihat{\wcns{c}}_i(\tup{y}_i)$.  In the
resulting disjunction, every disjunct $i' \neq i$ contains
$\Ihat{\wcns{c}}_i(\tup{y}_i) \feq \Ihat{\wcns{c}}_{i'}(\tup{y}'_{i'})$,
refuted by discrimination.  The remaining disjunct yields witnesses
$\tup{y}'_i$ with $\Ihat{\wcns{c}}_i(\tup{y}_i) \feq
\Ihat{\wcns{c}}_i(\tup{y}'_i)$ and $\Ihat{c}(\dots) \feq
R_i[\tup{y}'_i]$; injectivity gives $y_{ij} \feq y'_{ij}$, and rewriting
$R_i[\tup{y}'_i]$ back along these equations gives the guarded equation.

For strictness, note that all axioms of $\{\mathrm{Disj}_c\} \cup
\mathrm{Str}_I$ hold in the modification $\ML'$ of $\ML$ that
reinterprets $\Ihat{c}$ arbitrarily at arguments outside
$\semr{I}{}$-constructor forms with junk components --- e.g.\ setting
$\Ihat{c}^{\ML'}(\cls{\wcns{c}_i(\tup{e})}) \defeq \cls{\bx}$ whenever
some non-recursive $\cls{e_j}$ lies outside its argument-type
interpretation --- because every axiom of that set constrains
$\Ihat{c}$ only under guards.  The unguarded equation $\mathrm{Eq}_c$
fails in $\ML'$ at such a point whenever $R_i$'s value there differs
from $\cls{\bx}$, which is the case already for $\mathsf{mypred}$
(where $R_2 = y$): take $y \defeq \cls{\csO\,\csO} \neq \cls{\bx}$
(\cref{cor:discrimination}).
\end{proof}

Thus nothing is lost logically by splitting --- the disjunctive content
is recoverable with one inversion inference --- and something is gained:
the unguarded equations are strictly stronger, yet still sound
(\cref{prop:split-valid}).  The remaining difference is operational.

\subsection{Clause shapes and prover mechanics}
\label{sec:casesplit-cnf}

Consider $\mathsf{mypred}$ (\cref{ex:compile}).  Clausifying
$\mathrm{Disj}_{\mathsf{mypred}}$ requires a Skolem function $f$ for the
existential and distribution over the disjunction, yielding six
three-literal non-Horn clauses:
\[
\small
\begin{array}{ll}
  \neg \nat(x) \vee x \feq \csO \vee x \feq \csS(f(x)) &
  \neg \nat(x) \vee F(x) \feq \csO \vee x \feq \csS(f(x)) \\
  \neg \nat(x) \vee x \feq \csO \vee \nat(f(x)) &
  \neg \nat(x) \vee F(x) \feq \csO \vee \nat(f(x)) \\
  \neg \nat(x) \vee x \feq \csO \vee F(x) \feq f(x) &
  \neg \nat(x) \vee F(x) \feq \csO \vee F(x) \feq f(x)
\end{array}
\]
(writing $F$ for $\Ihat{\mathsf{mypred}}$), with two positive equality
literals per clause and the function's value expressed only through the
Skolem term.  Nested matches compound this multiplicatively.  A
superposition prover can use such clauses only through case-splitting
search: no clause is an orientable rewrite rule.  The split form
clausifies to
\[
  F(\csO) \feq \csO, \qquad F(\csS(y)) \feq y :
\]
two unit equations, orientable left-to-right under any simplification
ordering with the standard precedence heuristics, hence usable by
Vampire~\cite{KovacsVoronkov2013} and E~\cite{Schulz2002} as
\emph{demodulators} --- the prover evaluates $\mathsf{mypred}$ inside
terms at simplification cost, with no search, no Skolem functions and no
case splits.  For SMT solvers (CVC4/CVC5~\cite{BarrettEtAl2011},
Z3~\cite{deMouraBjorner2008}) the equations are ideal trigger material
for E-matching, whereas the disjunctive clauses produce case splits per
instantiation.  For recursive definitions the effect compounds: the
$\mathrm{Eq}$-form of $\mathsf{add}$ (\cref{ex:compile}) is exactly the
rewrite system a human would write, and equations of this shape are
well established as among the most valuable premises in the practice of
HOL-based hammers, where definitional packages export precisely the
per-constructor \texttt{simps} equations
\cite{Krauss2010,BlanchetteEtAl2016qed}, and dependently-typed frontends
export equational lemmas rather than raw match trees
\cite{QianEtAl2025}.

The one capability the disjunctive axiom packages more cheaply is a
\emph{case split on the scrutinee}: an ATP that has derived $\nat(t)$
obtains ``$t \feq \csO$ or $t \feq \csS(y)$'' together with the
corresponding value of $F$ in a single clause set.  In the split form
this requires chaining the (always emitted) inversion axiom with the
equations --- one extra inference, at the cost of deriving the guard
atom.  The exchange is measurable empirically; the a priori analysis
favours the split form on equational and computational goals and is
neutral on inversion-style goals, since inversion axioms are emitted
either way.

\begin{remark}[Reconstruction signal]
Since each split axiom is attached to one constructor, an ATP proof
that uses it reveals not only which match was relevant but which
\emph{branch} --- strictly more information for the reconstruction
backend (which maps used axioms back to case analyses) than the single
disjunctive axiom's name.
\end{remark}
